%%
% 第三章 系统架构设计（~10页）← 重点章
%%

\chapter{系统架构设计}

本章详细阐述基于多智能体系统（Multi-Agent System, MAS）\cite{guo2024llm_agent_survey, wu2023autogen}的自动化构建系统设计。多智能体系统是近年来随着大语言模型技术快速发展而兴起的系统架构范式，其核心思想是将复杂任务分解为多个专用智能体协同完成。与传统的单体系统相比，多智能体系统具有职责分离、可扩展性强和容错性好等优势，特别适合数据构建这类涉及多个异构处理阶段（PDF下载、多模态提取、文本清洗、引用图构建等）的复杂流程。AMMAG的自动化构建系统正是采用了这一架构范式，通过五个专用智能体的协同工作，实现了从arXiv元数据到标准化多模态属性图Benchmark数据集的端到端自动化。本章首先介绍系统的总体架构和设计原则，然后依次描述五个专用智能体的功能设计、核心算法实现和性能优化策略。

\section{系统总体架构}

\subsection{设计目标与原则}

AMMAG自动化构建系统基于多智能体系统架构，其设计目标是实现从arXiv元数据到标准化多模态属性图数据集的端到端自动化流程。学术Benchmark数据集的构建涉及多个异构处理阶段——从大规模元数据过滤、PDF批量下载、多模态内容提取、文本规范化清洗，到引用关系的提取与匹配、最终的数据集格式化——这些阶段在数据格式、计算资源需求和处理策略上差异显著。传统的单体系统难以同时满足这些异构需求，而多智能体系统通过将每个处理阶段封装为一个专用智能体，使每个智能体可以针对其特定任务进行独立优化。具体而言，系统需要满足以下四项设计目标：（1）端到端自动化——从原始arXiv论文列表到可直接用于图学习实验的标准化数据集，无需人工干预；（2）多模态提取——同时提取论文PDF中的文本、图像和数学公式三种模态信息；（3）引用图构建——自动匹配论文间的引用关系，构建引用网络图结构；（4）可扩展性——支持新增学术领域、新模态类型和新处理工具的灵活扩展。

基于上述目标，系统遵循以下四项设计原则：

\textbf{智能体职责分离原则。}这是多智能体系统最核心的设计原则。将复杂的Benchmark数据构建流程分解为多个功能独立的智能体，每个智能体专注于单一数据处理阶段。例如，论文下载智能体只需关注PDF获取，多模态提取智能体只需关注内容提取，引用图构建智能体只需关注引用关系匹配。职责分离使得每个智能体可以独立开发、测试和优化，降低了系统的整体复杂度。当某个阶段需要升级（如替换更先进的公式识别模型）时，只需修改对应的智能体，不影响其他智能体的功能。

\textbf{可扩展原则。}通过工具注册机制（MCP工具管理器）将底层工具与上层智能体解耦。新增处理工具只需实现标准化的工具接口并注册到工具管理器中，无需修改智能体或编排器的代码。这种设计使得多智能体系统可以灵活地扩展新的处理能力——例如，未来新增表格结构提取功能只需注册一个新的工具接口。

\textbf{智能体间松耦合原则。}智能体之间通过标准化的数据接口（任务配置字典和上游数据字典）进行通信，而非直接调用。每个智能体接收任务配置和上游智能体的输出作为输入，产生自己的输出传递给下游。这种松耦合设计使得智能体可以被替换、重排或并行执行，而不影响其他智能体的功能。

\textbf{确定性原则。}数据清洗和引用匹配等关键步骤采用纯规则驱动的方法实现，不依赖大语言模型，确保相同的输入始终产生相同的输出，保证数据处理的可复现性。这一原则确保了多智能体系统在数据处理层面的可靠性。

\subsection{三层架构设计}

系统采用编排器-智能体-工具三层多智能体架构，如图\ref{fig:architecture}所示。这一架构借鉴了AutoGen\cite{wu2023autogen}等多智能体框架的设计思想，将复杂的Benchmark构建任务组织为"编排器协调——智能体执行——工具支撑"的层次化协同模式。

\begin{figure}[htbp]
  \centering
  \includegraphics[width=0.95\textwidth]{images/architecture.png}
  \caption{基于多智能体系统的三层总体架构图}
  \label{fig:architecture}
\end{figure}

\textbf{编排器层}（Orchestrator Layer）位于多智能体系统的最上层，充当中央协调器的角色，负责调度五个智能体的执行顺序、管理智能体间的数据流转和处理异常情况。在多智能体系统架构中，编排器的存在是必要的——五个智能体各自独立运行，需要一个全局协调者来确保它们按照正确的顺序执行并传递中间结果。系统提供两种编排策略：顺序编排器（\texttt{DataCrawlerOrchestrator}）按固定顺序依次调度各智能体，适用于小规模数据处理；并行编排器（\texttt{ParallelDataCrawlerOrchestrator}）支持多GPU并行提取和多进程并行清洗，通过配置参数控制GPU数量和CPU工作进程数，适用于大规模数据处理。

\textbf{智能体层}（Agent Layer）是多智能体系统的核心层，包含五个功能专一、接口统一的专用智能体。多智能体系统的核心优势在于每个智能体可以针对其特定任务进行深度优化：（1）\texttt{ArxivBatchDownloaderAgent}——从arXiv元数据中按领域过滤论文，优化了大规模JSONL文件的流式读取；（2）\texttt{PaperCollectorAgent}——下载PDF并调用多模态提取引擎，优化了多GPU并行批量处理；（3）\texttt{DataCleanerAgent}——执行六步清洗流水线，采用纯规则驱动确保确定性；（4）\texttt{ReferenceGraphBuilderAgent}——构建引用网络图，设计了倒排索引和三级匹配引擎；（5）\texttt{DatasetFormatterAgent}——格式化为标准化数据集。每个智能体继承自抽象基类\texttt{BaseAgent}，实现统一的\texttt{execute()}接口，接收任务配置和上游智能体的输出数据，产生自己的输出传递给下游智能体。这种统一的接口设计是多智能体系统实现灵活编排的基础——编排器无需了解每个智能体的内部实现，只需调用其\texttt{execute()}方法并传递数据即可。

\textbf{工具层}（Tool Layer）位于多智能体系统的最下层，由\texttt{MCPToolManager}统一管理，为智能体层提供标准化的工具支撑。工具管理器采用插件式注册模式，将底层工具（PDF下载、多模态提取、数据清洗、格式转换等）封装为标准化的工具接口。智能体通过\texttt{call\_mcp\_tool()}方法调用工具，实现了智能体与具体工具实现的解耦。这种设计使得多智能体系统中的智能体无需关心底层工具的实现细节，降低了智能体的开发复杂度。

\subsection{数据流全流程}

系统的数据流如图\ref{fig:pipeline}所示。在多智能体系统架构下，完整的数据处理流程被分解为五个阶段，每个阶段由一个专用智能体负责，智能体之间通过标准化的数据接口传递中间结果：

\begin{figure}[htbp]
  \centering
  \includegraphics[width=0.95\textwidth]{images/pipeline.png}
  \caption{数据构建Pipeline全流程图}
  \label{fig:pipeline}
\end{figure}

（1）\texttt{ArxivBatchDownloaderAgent}从arXiv元数据集（约299万篇论文的JSON Lines文件）中按目标领域（如cs.SI）过滤论文，按更新日期排序后输出论文列表。

（2）\texttt{PaperCollectorAgent}接收论文列表，批量下载PDF文件，然后调用多模态提取引擎从每篇PDF中提取文本（Markdown格式）、图像（PNG格式）和公式（LaTeX格式）。

（3）\texttt{DataCleanerAgent}对提取的Markdown文本执行六步清洗流水线，修复PDF转换引入的格式问题，提取和格式化参考文献。

（4）\texttt{ReferenceGraphBuilderAgent}从PDF和Markdown双源提取引用信息，通过三级匹配引擎将引用与arXiv元数据中的论文进行匹配，构建引用网络图。

（5）\texttt{DatasetFormatterAgent}将清洗后的多模态数据、引用边信息和论文元数据整合为标准化的多模态属性图数据集。

\section{ArxivBatchDownloaderAgent}

ArxivBatchDownloaderAgent负责从大规模arXiv元数据集中按学术领域筛选目标论文。该智能体的核心功能是从约299万篇论文的元数据中高效地过滤出指定领域（如cs.SI）的论文，并按更新日期排序后输出给下游智能体。

\subsection{元数据过滤算法}

arXiv元数据集以JSON Lines格式存储，每行包含一篇论文的元数据信息（标题、作者、领域标签、摘要、更新日期等）。由于文件总大小约5.3GB，无法一次性加载到内存中。因此，系统采用流式读取策略，逐行解析JSON记录并按条件过滤。

过滤算法的核心逻辑为：对每条元数据记录，提取其\texttt{categories}字段的第一个标签作为主领域标签（primary category），判断是否与目标领域匹配。例如，一篇论文的\texttt{categories}字段为\texttt{"cs.SI cs.LG"}，则其主领域标签为\texttt{cs.SI}。通过仅匹配主领域标签，系统确保了筛选结果的领域纯度，避免了跨领域论文的引入。

匹配成功的论文按\texttt{update\_date}字段降序排序，优先选择最新更新的论文，然后取前\texttt{limit}篇作为最终结果。以cs.SI领域为例，系统从元数据中筛选出12,000篇候选论文，最终成功下载11,951篇（下载成功率99.6\%）。

\section{PaperCollectorAgent与多模态提取引擎}

PaperCollectorAgent是系统中计算密集度最高的智能体，负责批量下载论文PDF并调用多模态提取引擎从中提取文本、图像和公式三种模态信息。

\subsection{多模态提取引擎}

多模态提取引擎基于PDF-Extract-Kit工具包\cite{zhao2024doclayoutyolo}实现，其处理流程如图\ref{fig:extraction}所示。对于每篇论文PDF，引擎首先使用DocLayout-YOLO\cite{zhao2024doclayoutyolo}进行版面检测，识别出页面中的五种区域类型：文本块、图片、行内公式、行间公式和表格。

\begin{figure}[htbp]
  \centering
  \includegraphics[width=0.95\textwidth]{images/extraction.png}
  \caption{多模态提取引擎处理流程}
  \label{fig:extraction}
\end{figure}

版面检测完成后，引擎根据区域类型分别进行处理：

\textbf{文本区域：}使用PaddleOCR\cite{du2020ppocr}对文本区域进行光学字符识别（OCR），将识别结果按段落结构组织为Markdown格式。PaddleOCR是一个开源的超轻量级OCR系统，支持多语言文字识别，在中文和英文文档上均有良好的识别精度。

\textbf{图片区域：}直接从PDF中裁剪图片区域的像素数据，保存为PNG格式。每篇论文的所有图片存储在独立的\texttt{images/}子目录中。

\textbf{公式区域：}首先使用YOLOv8检测行内和行间公式区域，然后使用UniMERNet\cite{wang2024unimernet}将公式图像识别为LaTeX代码。UniMERNet是一个通用的数学公式识别网络，能够处理复杂的手写和印刷体数学表达式，包括分数、根号、矩阵、积分等高级数学结构。

每篇论文的提取结果以统一格式存储：\texttt{paper.md}（Markdown全文）、\texttt{paper.npy}（结构化数据文件）和\texttt{images/}（图片目录）。

\subsection{多GPU并行提取}

多模态提取是整个系统中计算开销最大的环节。每次启动提取模型（DocLayout-YOLO、UniMERNet等）需要约30秒的模型加载时间。如果逐篇处理，仅模型加载就需消耗大量时间。

系统设计了\texttt{GPUExtractionPool}实现多GPU并行批量提取。其核心优化策略如下：

\textbf{批量处理。}将多篇论文的PDF通过符号链接组织到临时目录中，使用\texttt{-{}-pdf\_dir}参数一次性提交给提取引擎。这样，模型只需加载一次即可处理整个批次的论文（默认批次大小为50篇），将每篇论文的模型加载开销从30秒降至约0.6秒。

\textbf{多GPU并行。}系统将待处理的论文列表按轮转（round-robin）方式分配到多个GPU上，每个GPU由独立的线程驱动。默认配置使用8个GPU，通过\texttt{CUDA\_VISIBLE\_DEVICES}环境变量隔离每个GPU的可见设备，确保不同GPU上的提取进程互不干扰。

\textbf{失败回退。}批量提取中如果某篇PDF处理失败（如损坏的PDF文件），不会阻塞整个批次。系统对失败的论文回退到逐篇提取模式（fallback sequential），单独处理每篇论文，保证单篇论文的失败不影响批次中其他论文的处理。

\textbf{跳过已处理。}在开始提取前，系统检查每篇论文的输出目录是否已存在\texttt{.md}或\texttt{.npy}文件。如果已存在，则跳过该论文的提取，避免重复计算。

\section{DataCleanerAgent}

DataCleanerAgent负责对提取的Markdown文本进行规则驱动的清洗处理。该智能体不使用大语言模型，完全通过正则表达式和启发式规则实现，确保清洗过程的确定性和可复现性。

\subsection{清洗流水线}

清洗流水线由六个步骤组成，依次对Markdown文本进行规范化处理：

\textbf{步骤一：断词修复。}PDF转换过程中，英文单词经常被连字符和换行符分割。例如，"learning"可能被转换为"learn-\\ning"。断词修复步骤使用正则表达式\texttt{r"(\textbackslash w+)-\textbackslash s*\textbackslash n\textbackslash s*(\textbackslash w+)"}匹配被分割的单词模式，将断词部分重新合并为完整的单词。

\textbf{步骤二：段落合并。}PDF提取的文本经常出现段落被错误分行的情况。段落合并步骤逐行扫描文本，判断相邻两行是否应合并为同一段落。合并的条件包括：当前行不以句号、感叹号、问号等句末标点结尾；下一行不以Markdown标题（\texttt{\#}）、列表标记（\texttt{-}、\texttt{*}、\texttt{1.}）或引用编号（\texttt{[N]}）开头；下一行不以大写字母+逗号+首字母缩写模式开头（避免合并参考文献的作者列表）。

\textbf{步骤三：重复段落去除。}使用MD5哈希值对段落进行去重。对于长度超过50个字符的段落，计算其去除空白字符后的MD5哈希值，如果该哈希值已被记录则跳过该段落。长度不足50字符的短段落无条件保留，避免误删标题等短文本。

\textbf{步骤四：公式格式标准化。}统一LaTeX公式的格式规范：去除\texttt{\$\$}前后多余的空白字符；确保块级公式（\texttt{\$\$}）独占一行；统一行内公式（\texttt{\$...\$}）内部的前后空格。

\textbf{步骤五：参考文献提取与格式化。}这是最复杂的清洗步骤。系统使用三种策略依次尝试提取和格式化参考文献：（1）按\texttt{[N]}数字标签分割参考文献条目；（2）按作者姓名模式分割；（3）按年份边界分割。每种策略提取的参考文献条目都需通过有效性验证（包含4位年份、长度至少30字符、包含作者姓名模式或已知会议/期刊名称）。最终将参考文献按序编号为\texttt{[1], [2], ...}的标准格式。

\textbf{步骤六：冗余空白清理。}压缩三个及以上连续空行为两个空行，去除每行末尾的空白字符，确保文件以单个换行符结尾。

\subsection{并行清洗}

清洗步骤支持多进程并行执行。系统使用Python的\texttt{multiprocessing.Pool}创建指定数量的工作进程（默认8个），每个进程独立执行完整的六步清洗流水线。由于清洗过程是纯CPU操作且无共享状态，多进程并行可以实现近线性的加速比。

\section{ReferenceGraphBuilderAgent}

ReferenceGraphBuilderAgent是系统中技术复杂度最高的智能体，负责从论文中提取引用信息并构建引用网络图。该智能体采用双源提取策略获取原始引用信息，通过三级匹配引擎将引用与arXiv元数据中的论文进行匹配，最终输出引用边的CSV文件。

\subsection{双源引用提取}

学术论文的参考文献信息可以从两个数据源提取，两个源具有互补性：

\textbf{数据源一：PDF直接解析。}使用pdfplumber库直接解析PDF文件，提取参考文献页面的原始文本。该数据源的优势在于保留了PDF的精确排版信息，包括双栏布局。系统通过分析页面中字符的x坐标分布来检测双栏布局——在页面中间40\%的宽度范围内寻找最大的x坐标间隙，如果该间隙超过12磅（point），则判定为双栏布局，分别裁剪左右两栏进行文本提取。此外，系统通过检测跨页面重复的首行和末行来识别并移除页眉和页脚。

\textbf{数据源二：Markdown提取。}从已提取的Markdown文本中搜索参考文献章节。系统使用五种正则表达式模式匹配不同的章节标题格式（\texttt{\#\# References}、\texttt{\# References}、\texttt{References\\n}等）。对每个候选区域，系统计算一个综合评分函数，根据引用编号密度、arXiv标识符密度、年份密度等指标评估该区域作为参考文献章节的可能性。

两种数据源的融合显著提升了引用信息的完整性。在小规模验证实验中，对于cs.SI领域的论文样本，单源提取仅获得344条引用边，而双源融合后获得1,064条引用边，提升了209\%。

\subsection{多格式引用解析}

学术论文的参考文献采用多种不同的格式，系统支持识别和解析七种主要的引用格式类型：（1）方括号数字编号\texttt{[N]}格式；（2）圆括号数字编号\texttt{(N)}格式；（3）作者-年份编码\texttt{[Author+YY]}格式；（4）全角方括号数字编号格式；（5）点号编号\texttt{N.}格式；（6）空格编号\texttt{N Author}格式；（7）无编号格式（通过行级别或年份边界分割）。

PDF提取过程中，多条引用记录可能被错误地合并为一行（特别是当引用条目较短时）。系统设计了多信号边界检测算法来处理这种情况。该算法为四种边界信号设定权重：方括号编号（权重10）、点号编号（权重8，后跟大写字母时增强到10）、年份边界（权重3）和arXiv标识符边界（权重5）。相邻10个字符内的边界信号被合并（保留较高权重），最终需要至少2个有效边界（权重$\geq$3）才执行分割。

此外，系统对分割后的片段进行分类（arXiv尾部、DOI尾部、出版商名称、年份尾部、URL尾部、OCR碎片等），将短片段合并到前一条引用记录中，避免信息丢失。

\subsection{三级匹配引擎}

三级匹配引擎是引用图构建的核心算法，其设计目标是高效、准确地将引用记录中的论文标题与arXiv元数据中的论文进行匹配。

\textbf{Level 0：arXiv标识符直接匹配。}如果引用记录中包含arXiv标识符（格式为\texttt{arXiv:NNNN.NNNNN}或\texttt{arxiv.org/abs/NNNN.NNNNN}），系统直接在标识符映射表中进行O(1)哈希查找。这是最可靠的匹配方式，因为arXiv标识符具有全局唯一性。

\textbf{Level 1：精确标题匹配。}将引用记录中提取的论文标题进行归一化处理（转小写、去除标点、统一空白字符），然后在标题-标识符映射表中进行精确查找。归一化处理消除了大小写和标点差异对匹配的影响。

\textbf{Level 2：模糊标题匹配。}当精确匹配失败时，系统使用基于倒排索引\cite{zobel2006inverted, manning2008ir}的模糊匹配算法。该算法分为三个阶段：

（a）\textit{候选生成。}将归一化后的引用标题分词，去除停用词（31个常见英文停用词）和长度小于3的词，保留"显著词"。对每个显著词，通过倒排索引查找包含该词的所有元数据标题，统计每个候选标题与查询标题的显著词重叠数。仅保留重叠数$\geq$2的候选标题。

（b）\textit{候选过滤。}对候选标题进行长度过滤——如果候选标题的归一化长度与查询标题相差超过50\%，则丢弃该候选。剩余候选按重叠数降序排列，仅评估前500个。

（c）\textit{匹配验证。}对每个候选标题，依次尝试两种匹配策略：

\hspace{1em}\textbullet~\textit{双向子串匹配}：检查两个标题是否存在子串包含关系（长度比率$>$0.75），并验证关键词重叠率$\geq$0.35。

\hspace{1em}\textbullet~\textit{SequenceMatcher模糊匹配}\cite{ratcliff1988gestalt}：使用Python \texttt{difflib.SequenceMatcher}计算两个标题字符串的相似度比率。设置匹配阈值为0.90（即两个标题至少有90\%的字符匹配），并要求关键词重叠率$\geq$0.30。当相似度达到0.95时，直接返回匹配结果（早期退出）。

所有匹配结果都经过arXiv标识符交叉验证：当通过标题匹配到一篇论文时，如果引用记录中也包含arXiv标识符，系统验证匹配论文的标识符与引用记录中的标识符是否一致（关键词重叠率$\geq$0.40），防止误匹配。

以cs.SI领域为例，系统从11,929篇论文中提取了21,638条原始引用边，经过三级匹配引擎处理后，成功匹配4,483条域内引用边，匹配率为20.7\%。未匹配的引用主要指向域外论文（其他领域的论文），这一匹配率在学术引用图的构建中是合理的。

\subsection{并行化与零拷贝优化}

引用图构建需要访问约299万篇论文的元数据索引（标题-标识符映射表、标识符-标题映射表、倒排索引等），这些索引数据结构的总内存占用约为数GB。如果使用Python的\texttt{multiprocessing.Pool}在多个工作进程间共享这些数据，传统的序列化-反序列化方式将产生巨大的开销。

系统采用fork-based Copy-on-Write（COW）零拷贝优化来解决这个问题。在创建进程池之前，主进程加载全部元数据并构建索引数据结构，将其存储为模块级别的全局变量。当\texttt{multiprocessing.Pool}使用\texttt{fork()}系统调用创建子进程时，子进程通过操作系统的Copy-on-Write机制继承父进程的内存空间。由于子进程只读取这些数据而不修改它们，父进程和所有子进程共享同一份物理内存页，实现了真正的零拷贝共享。这种优化避免了GB级别的数据序列化开销，显著提升了并行引用图构建的效率。

\section{并行化与性能优化}

本节总结系统中的并行策略及其性能优化效果。

\subsection{多GPU并行提取}

多模态提取是整个系统中计算开销最大的环节。系统设计了\texttt{GPUExtractionPool}实现多GPU并行批量提取（详见3.4.2节），将待处理的论文列表按轮转方式分配到多个GPU上，每个GPU由独立线程驱动，并通过批量处理策略将模型加载开销从每篇30秒降至每批30秒。

\subsection{多进程并行清洗}

数据清洗阶段采用\texttt{multiprocessing.Pool}创建指定数量的CPU工作进程（默认32个），每个进程独立执行完整的六步清洗流水线。由于清洗过程是纯CPU操作且无共享状态，多进程并行可以实现近线性的加速比。

\subsection{COW零拷贝优化}

引用图构建需要访问约299万篇论文的元数据索引（标题-标识符映射表、倒排索引等），总内存占用约为数GB。系统采用fork-based Copy-on-Write（COW）零拷贝优化来避免GB级别数据的序列化开销。在创建进程池之前，主进程加载全部元数据并构建索引数据结构，将其存储为模块级别的全局变量。当\texttt{multiprocessing.Pool}使用\texttt{fork()}系统调用创建子进程时，子进程通过操作系统的Copy-on-Write机制继承父进程的内存空间。由于子进程只读取这些数据而不修改它们，父进程和所有子进程共享同一份物理内存页，实现了真正的零拷贝共享。

\subsection{性能对比}

表\ref{tab:performance}总结了不同配置下的处理性能。以cs.SI领域（11,929篇论文）为例，从元数据过滤到最终数据集输出的完整流程，在8 GPU并行提取和32 CPU并行清洗的配置下，相比串行模式实现了约7倍的总处理时间缩减。

\begin{table}[htbp]
  \centering
  \caption{不同配置下的处理性能（cs.SI, 11,929篇论文）}
  \label{tab:performance}
  \begin{tabular}{lccc}
    \toprule
    配置 & 模型加载 & 提取并行度 & 估算总时间 \\
    \midrule
    串行（1 GPU, 1 CPU） & 每篇30s & 1 & $\sim$99h \\
    多GPU批量（8 GPU） & 每批30s & 8 & $\sim$14h \\
    \bottomrule
  \end{tabular}
\end{table}

\section{MCP工具管理层}

MCP工具管理层是系统三层架构的基础层，负责管理和调度所有底层工具。工具管理器采用插件式注册模式，支持四类共13个工具接口。

\textbf{论文下载工具组：}包括单篇下载工具（支持重试和指数退避）、批量下载工具（使用线程池并行下载，最大10个线程）和DOI下载工具。下载工具在下载前移除所有代理环境变量，避免代理配置导致的连接问题。

\textbf{多模态提取工具组：}包括单篇提取工具、批量提取工具以及文本/图片/表格提取工具。核心提取工具封装了PDF-Extract-Kit的子进程调用，设置30分钟超时防止处理卡死。

\textbf{数据清洗工具组：}包括文本清洗、数据去重和数据验证三个工具接口。

\textbf{数据集工具组：}包括属性图格式化和数据集保存两个工具接口。

工具管理器的核心接口为\texttt{call\_tool(name, **kwargs)}，智能体通过工具名称和参数调用对应的工具函数。这种设计使得新增工具只需实现函数签名并注册到管理器中，无需修改智能体或编排器的代码，实现了工具层与智能体层的完全解耦。这种设计借鉴了AutoGen\cite{wu2023autogen}等框架的工具管理思想，通过标准化的工具接口实现工具的即插即用——其设计理念与模型上下文协议（Model Context Protocol, MCP）类似。

\section{本章小结}

本章详细阐述了基于多智能体系统（MAS）的自动化构建系统的三层架构设计。系统采用编排器-智能体-工具三层多智能体架构，包含五个专用智能体，通过编排器的中央协调实现了从arXiv元数据到标准化多模态属性图数据集的端到端自动化流程。多智能体系统架构的引入是本系统设计的核心决策——学术Benchmark数据集的构建涉及多个异构处理阶段，多智能体系统通过职责分离使每个智能体专注于单一任务，通过松耦合的接口设计实现灵活编排，通过独立的性能优化（如多GPU并行提取、多进程并行清洗、COW零拷贝）使每个智能体在其特定任务上达到最优性能。这种架构模式使得系统具备良好的可扩展性——未来新增处理阶段（如表格结构提取、LLM辅助引用匹配等）只需实现新的智能体并注册到编排器中。

系统的核心技术创新包括：（1）基于DocLayout-YOLO、UniMERNet和PaddleOCR的多模态提取引擎，实现了论文PDF中文本、图像和公式的并行提取；（2）双源引用提取策略，通过PDF直接解析和Markdown提取两种互补数据源，将引用信息的完整性提升了209\%；（3）三级匹配引擎，结合arXiv标识符直接匹配、精确标题匹配和基于倒排索引的模糊匹配，实现了高效的引用-论文匹配；（4）多GPU并行提取和多进程并行清洗，通过批量处理和进程池策略实现了高效的数据处理；（5）fork-based COW零拷贝优化，避免了GB级别元数据索引的序列化开销。
